\documentclass[a4paper,12pt]{article}
\usepackage{docsTemplate}

\setTitle{Manuale Utente}
\setAuthors{}
\setVerificators{}
\setApprovation{}
\setVersion{0.0}
\setType{Esterno}
\setDestination{Team Three Technologies \\ & Prof. Tullio Vardanega \\ & Prof. Riccardo Cardin \\ & Sanmarco Informatica SpA}

\begin{document}
    \include{docTitlepage}

    \section*{Registro delle modifiche} {
        \begin{table}[h!]
            \rowcolors{2}{lightgray}{white}
                \begin{tabularx}{\textwidth}{|l|l|l|l|L|L|}
                \hline
                \textbf{Vers.} & \textbf{Data} & \textbf{Autore} & \textbf{Ruolo} & \textbf{Verifica} & \textbf{Oggetto} \\
                \hline
            \end{tabularx}
        \end{table}
    }
    \newpage
    
    \tableofcontents
    \newpage

    \listoffigures
    \newpage
    
    \listoftables
    \newpage

    \section{Introduzione} {
        \subsection{Scopo del documento} {        
            Il presente documento ha lo scopo di fornire all'utente tutte le istruzioni necessarie per la corretta installazione e il normale utilizzo dell'applicativo desktop DIPReader.
            Al suo interno è presente una guida per la consultazione dei pacchetti informativi (DIP).
        }

        \subsection{Glossario} {
            Si rimanda al documento Glossario per chiarire i termini nel dominio del progetto che possono risultare ambigui o di difficile comprensione. I vocaboli presenti nel glossario sono stati segnalati seguendo la notazione:
            \begin{center}
                parola\textsubscript{G}
            \end{center}
        }
        
        \subsection{Riferimenti} {
            \subsubsection{Riferimenti normativi} {
                \begin{itemize}
                    \item
                \end{itemize}
            }
            \subsubsection{Riferimenti informativi} {
                \begin{itemize}
                    \item \textit{Glossario v0.5}
                    \begin{itemize}
                        \item \url{https://team-three-technologies.github.io/Documenti/1%20-%20RTB/Documenti%20interni/Glossario_v0_5.pdf}
                        \item[] (Ultimo accesso: 2026/xx/xx)
                    \end{itemize}
                \end{itemize}
            }
        }
    }
    
    \section{Installazione}
        \subsection{Requisiti di sistema}

        DIPReader è un'applicazione desktop autonoma che non richiede l'installazione di software aggiuntivo. Tutto il necessario per il funzionamento è già incluso nel pacchetto di distribuzione. 
        L'applicazione è compatibile con i seguenti sistemi operativi:
        \begin{itemize}
            \item Windows: Windows 10 o versioni successive
            \item macOS: macOS 11 o versioni successive
            \item Linux: qualsiasi distribuzione moderna con supporto al formato AppImage
        \end{itemize}

        \subsection{Installazione}
            come installare l'applicativo

        \subsection{Primo avvio} {
            ...
        }

    \section{Interfaccia}
        panoramica dell'Interfaccia

    \section{Guida all'utilizzo}
        \subsection{Apertura di un pacchetto DIP}

        \subsection{Consultazione dei documenti}
            \subsubsection{Visualizzare la lista dei documenti}
            \subsubsection{Selezionare un documento e visualizzarne i dettagli}
            \subsubsection{Visualizzare l'anteprima di un documento}

        \subsection{Consultazione degli allegati}
            \subsubsection{Visualizzare gli allegati di un documento}
            \subsubsection{Selezionare un allegato e visualizzarne i dettagli}
            \subsubsection{Tornare al documento padre}
        
        \subsection{Ricerca e uso dei filtri}
            \subsubsection{Aggiungere un filtro}
            \subsubsection{Utilizzare più filtri combinati}
            \subsubsection{Rimuovere un filtro}
            \subsubsection{Avviare la ricerca}
            
        \subsection{Visualizzare i dettagli del DIP}

\end{document}